\scanpage{19}
% P01: the first dotted T_i has a heavily retraced, forked body; no earlier glyph or cancellation state is established.
mieux, si la \uline{famille} des $\dot T_i\subset\dot T$ au dessus
de $Y$ est \deleted{\unsure{globalement}} loc. triviale. On étend cette
définition au cas de $X_i\subset X$, et d'un
voisinage \deleted{tubulaire}\inserted{cônique} $T$ de $Y$ dans $X$ : on dit
il est équisingulier p.r. à $(X_i)$, si $T$
% P04: the detached low point after the first T is retained literally; an i-index remains an alternative, with no assigned mathematical role.
l'est par rapport à $(X_i\cap\unsure{T.|T_i}\subset T)_i$. Cela
implique que si $\forall i$, $X_i\cap Y$ est une
partie à la fois ouverte et fermée de $Y$
alors, si $Y$ est connexe, que \deleted{$X_i\cap Y$} l'on ait
l'alternative
\[
% P07: the T roof has a retraced upper loop, without an independent overbar.
(14)\qquad X_i\cap T=\phi\quad\underline{\text{ou}}\quad X_i\supset Y\ .
\]
\margin{on dit qu'il est \uline{strictement équisingulier}
si de plus il est \uline{strictement} \unsure{compatible} avec tous
les $X_i$.}
\margin{Si $T$ équisingulier p.r. à $(X_i)$,
$\lambda T$ l'est strictement pour $\lambda\in]0,1[$.}

% P10: the cancellation continues above the circled number in a dense loop; no separate additional digit is recoverable.
\deleted{\textcircled{\unsure{5|8}}}
On dit aussi que $X,(X_i)_{i\in I}$ est
\inserted{(globalement)} équisingulier
le long de $Y$, s'il existe un voisinage cônique
% P12: the point above the retraced T is separately written and live; the literal dotted form is retained.
\deleted{\unsure{de}} $\dot T$ de $Y$ dans $X$, qui est équisingulier
p.r. \deleted{\unsure{à}} $(X_i)$ --- auquel cas les $\lambda T$
$(\lambda\in]0,1[)$ sont $=$ strictement équisinguliers.
\margin{18}
